\documentclass[10pt,a4paper]{amsart}
\usepackage{amsmath,amssymb,amsthm,graphicx}
% Retained preamble for the Stacks Project chapter snapshot in this directory.
% The chapter's own preamble.tex is a copy of the shared Stacks Project preamble; the block
% below is sliced MECHANICALLY, line for line, from it. Nothing was added, removed or
% redefined; the wrapper class, the xr-hyper/hyperref wiring and the externaldocument list
% are omitted because the assembler supplies the document class.
%
% lines taken from preamble.tex: 3, 5, 9, 11, 12, 156, 157, 158, ... (51 lines total)
%
% For commutative diagrams we use Xy-pic
\usepackage[all]{xy}

% We use 2cell for 2-commutative diagrams.
\xyoption{2cell}
\UseAllTwocells

\newenvironment{reference}{\comment}{\endcomment}
\newenvironment{slogan}{\comment}{\endcomment}
\newenvironment{history}{\comment}{\endcomment}
\theoremstyle{plain}
\newtheorem{theorem}[subsection]{Theorem}
\newtheorem{proposition}[subsection]{Proposition}
\newtheorem{lemma}[subsection]{Lemma}
\theoremstyle{definition}
\newtheorem{definition}[subsection]{Definition}
\newtheorem{example}[subsection]{Example}
\newtheorem{exercise}[subsection]{Exercise}
\newtheorem{situation}[subsection]{Situation}
\theoremstyle{remark}
\newtheorem{remark}[subsection]{Remark}
\newtheorem{remarks}[subsection]{Remarks}
\numberwithin{equation}{subsection}
\def\lim{\mathop{\mathrm{lim}}\nolimits}
\def\colim{\mathop{\mathrm{colim}}\nolimits}
\def\Spec{\mathop{\mathrm{Spec}}}
\def\Hom{\mathop{\mathrm{Hom}}\nolimits}
\def\Ext{\mathop{\mathrm{Ext}}\nolimits}
\def\SheafHom{\mathop{\mathcal{H}\!\mathit{om}}\nolimits}
\def\SheafExt{\mathop{\mathcal{E}\!\mathit{xt}}\nolimits}
\def\Sch{\mathit{Sch}}
\def\Mor{\mathop{\mathrm{Mor}}\nolimits}
\def\Ob{\mathop{\mathrm{Ob}}\nolimits}
\def\Sh{\mathop{\mathit{Sh}}\nolimits}
\def\NL{\mathop{N\!L}\nolimits}
\def\CH{\mathop{\mathrm{CH}}\nolimits}
\def\proetale{{pro\text{-}\acute{e}tale}}
\def\etale{{\acute{e}tale}}
\def\QCoh{\mathit{QCoh}}
\def\Ker{\mathop{\mathrm{Ker}}}
\def\Im{\mathop{\mathrm{Im}}}
\def\Coker{\mathop{\mathrm{Coker}}}
\def\Coim{\mathop{\mathrm{Coim}}}
\DeclareMathSymbol{\boxtimes}{\mathbin}{AMSa}{"02}
\def\QCohstack{\mathcal{QC}\!\mathit{oh}}
\def\Cohstack{\mathcal{C}\!\mathit{oh}}
\def\Spacesstack{\mathcal{S}\!\mathit{paces}}
\def\Quotfunctor{\mathrm{Quot}}
\def\Hilbfunctor{\mathrm{Hilb}}
\def\Curvesstack{\mathcal{C}\!\mathit{urves}}
\def\Polarizedstack{\mathcal{P}\!\mathit{olarized}}
\def\Complexesstack{\mathcal{C}\!\mathit{omplexes}}
\def\Pic{\mathop{\mathrm{Pic}}\nolimits}
\def\Picardstack{\mathcal{P}\!\mathit{ic}}
\def\Picardfunctor{\mathrm{Pic}}
\def\Deformationcategory{\mathcal{D}\!\mathit{ef}}
\usepackage{cleveref}
\crefname{originalreference}{original reference}{original references}
\Crefname{originalreference}{Original reference}{Original references}
\makeatletter
\newlabel{stacks-sheaves-lemma-lqc-flat-base-change-fppf-sheaf}{{076T}{source}{}{sourcepointer-1}{}}
\newlabel{stacks-sheaves-lemma-lqc-flat-base-change-fppf-sheaf@cref}{{[originalreference][0][]076T}{[1][0][]source}{}{}{}}
\makeatother
\crefname{originallemma}{original lemma}{original lemmas}
\Crefname{originallemma}{Original lemma}{Original lemmas}
\makeatletter
\newlabel{lemma-flat-comparison}{{0765}{source}{}{sourcepointer-2}{}}
\newlabel{lemma-flat-comparison@cref}{{[originallemma][0][]0765}{[1][0][]source}{}{}{}}
\makeatother
\crefname{originalreference}{original reference}{original references}
\Crefname{originalreference}{Original reference}{Original references}
\makeatletter
\newlabel{stacks-sheaves-lemma-compare-fppf-etale}{{075R}{source}{}{sourcepointer-3}{}}
\newlabel{stacks-sheaves-lemma-compare-fppf-etale@cref}{{[originalreference][0][]075R}{[1][0][]source}{}{}{}}
\makeatother
\crefname{originalreference}{original reference}{original references}
\Crefname{originalreference}{Original reference}{Original references}
\makeatletter
\newlabel{sites-cohomology-lemma-kill-cohomology-class-on-covering}{{01FW}{source}{}{sourcepointer-4}{}}
\newlabel{sites-cohomology-lemma-kill-cohomology-class-on-covering@cref}{{[originalreference][0][]01FW}{[1][0][]source}{}{}{}}
\makeatother
\crefname{originallemma}{original lemma}{original lemmas}
\Crefname{originallemma}{Original lemma}{Original lemmas}
\makeatletter
\newlabel{lemma-check-flat-comparison-on-etale-covering}{{0764}{source}{}{sourcepointer-5}{}}
\newlabel{lemma-check-flat-comparison-on-etale-covering@cref}{{[originallemma][0][]0764}{[1][0][]source}{}{}{}}
\makeatother
\crefname{originalreference}{original reference}{original references}
\Crefname{originalreference}{Original reference}{Original references}
\makeatletter
\newlabel{stacks-sheaves-lemma-pullback-lqc}{{06WL}{source}{}{sourcepointer-6}{}}
\newlabel{stacks-sheaves-lemma-pullback-lqc@cref}{{[originalreference][0][]06WL}{[1][0][]source}{}{}{}}
\makeatother
\crefname{originalreference}{original reference}{original references}
\Crefname{originalreference}{Original reference}{Original references}
\makeatletter
\newlabel{stacks-sheaves-lemma-lqc-colimits}{{06WM}{source}{}{sourcepointer-7}{}}
\newlabel{stacks-sheaves-lemma-lqc-colimits@cref}{{[originalreference][0][]06WM}{[1][0][]source}{}{}{}}
\makeatother
\title{Summary of results on locally quasi-coherent modules with the flat base change property}
\author{The Stacks Project Authors}
\date{stacks-project:stacks-cohomology.tex:local-snapshot-20261012; GNU FDL 1.2 or later}
\begin{document}
\maketitle
\noindent\textit{Editorial note.} All delivered mathematical text is the author's own, copied verbatim from the pinned Stacks Project chapter file stacks-cohomology.tex; no statement or proof step was rewritten, reordered or removed.\par\medskip
\noindent\textit{Allowed standard prerequisites (editorial).}\begin{itemize}
\item fppf-etale cohomology comparison and locality of cohomology. Sheaves on Stacks, Lemma stacks-sheaves-lemma-compare-fppf-etale; Cohomology on Sites, Lemma sites-cohomology-lemma-kill-cohomology-class-on-covering. Part (2).
\item delivered companions: flat comparison and pushforward stability. delivered companions lemma-flat-comparison, lemma-pushforward-locally-quasi-coherent, lemma-check-lqc-fbc-on-covering. Parts (3)-(7).
\end{itemize}\medskip
\ifdefined\hypertarget\hypertarget{sourcepointer-1}{}\fi
\noindent\textit{Source reference.} Original reference 076T: sources/stacks-chow/tags.txt, source lines 9178--9178. Cross-chapter standard prerequisite cited by the author inside the counted span; it is used as a pointer and its own proof is not part of this item.\par\medskip
\ifdefined\hypertarget\hypertarget{sourcepointer-2}{}\fi
\noindent\textit{Source reference.} Original lemma 0765: sources/stacks-chow/tags.txt, source lines 9155--9155. Cross-chapter standard prerequisite cited by the author inside the counted span; it is used as a pointer and its own proof is not part of this item.\par\medskip
\ifdefined\hypertarget\hypertarget{sourcepointer-3}{}\fi
\noindent\textit{Source reference.} Original reference 075R: sources/stacks-chow/tags.txt, source lines 9140--9140. Cross-chapter standard prerequisite cited by the author inside the counted span; it is used as a pointer and its own proof is not part of this item.\par\medskip
\ifdefined\hypertarget\hypertarget{sourcepointer-4}{}\fi
\noindent\textit{Source reference.} Original reference 01FW: sources/stacks-chow/tags.txt, source lines 1885--1885. Cross-chapter standard prerequisite cited by the author inside the counted span; it is used as a pointer and its own proof is not part of this item.\par\medskip
\ifdefined\hypertarget\hypertarget{sourcepointer-5}{}\fi
\noindent\textit{Source reference.} Original lemma 0764: sources/stacks-chow/tags.txt, source lines 9154--9154. Cross-chapter standard prerequisite cited by the author inside the counted span; it is used as a pointer and its own proof is not part of this item.\par\medskip
\ifdefined\hypertarget\hypertarget{sourcepointer-6}{}\fi
\noindent\textit{Source reference.} Original reference 06WL: sources/stacks-chow/tags.txt, source lines 8820--8820. Cross-chapter standard prerequisite cited by the author inside the counted span; it is used as a pointer and its own proof is not part of this item.\par\medskip
\ifdefined\hypertarget\hypertarget{sourcepointer-7}{}\fi
\noindent\textit{Source reference.} Original reference 06WM: sources/stacks-chow/tags.txt, source lines 8821--8821. Cross-chapter standard prerequisite cited by the author inside the counted span; it is used as a pointer and its own proof is not part of this item.\par\medskip
\begin{proposition}
\label{proposition-loc-qcoh-flat-base-change}
Summary of results on locally quasi-coherent modules having the flat
base change property.
\begin{enumerate}
\item Let $\mathcal{X}$ be an algebraic stack.
If $\mathcal{F}$ is in $\textit{LQCoh}^{fbc}(\mathcal{O}_\mathcal{X})$,
then $\mathcal{F}$ is a sheaf for the fppf topology, i.e., it is
an object of $\textit{Mod}(\mathcal{O}_\mathcal{X})$.
\item The category $\textit{LQCoh}^{fbc}(\mathcal{O}_\mathcal{X})$
is a weak Serre subcategory of both $\textit{Mod}(\mathcal{O}_\mathcal{X})$
and $\textit{Mod}(\mathcal{X}_\etale, \mathcal{O}_\mathcal{X})$.
\item Pullback $f^*$ along any morphism of algebraic stacks
$f : \mathcal{X} \to \mathcal{Y}$ induces a functor
$f^* : \textit{LQCoh}^{fbc}(\mathcal{O}_\mathcal{Y}) \to
\textit{LQCoh}^{fbc}(\mathcal{O}_\mathcal{X})$.
\item If $f : \mathcal{X} \to \mathcal{Y}$ is a
quasi-compact and quasi-separated morphism of algebraic stacks
and $\mathcal{F}$ is an object of
$\textit{LQCoh}^{fbc}(\mathcal{O}_\mathcal{X})$, then
\begin{enumerate}
\item the total direct image $Rf_*\mathcal{F}$ and the higher direct
images $R^if_*\mathcal{F}$ can be computed in either the \'etale or the
fppf topology with the same result, and
\item each $R^if_*\mathcal{F}$ is an object of
$\textit{LQCoh}^{fbc}(\mathcal{O}_\mathcal{Y})$.
\end{enumerate}
\item The category $\textit{LQCoh}^{fbc}(\mathcal{O}_\mathcal{X})$ has
colimits and they agree with colimits in
$\textit{Mod}(\mathcal{X}_\etale, \mathcal{O}_\mathcal{X})$
as well as in $\textit{Mod}(\mathcal{O}_\mathcal{X})$.
\item Given $\mathcal{F}$ and $\mathcal{G}$ in
$\textit{LQCoh}^{fbc}(\mathcal{O}_\mathcal{X})$ then the tensor product
$\mathcal{F} \otimes_{\mathcal{O}_\mathcal{X}} \mathcal{G}$
is in $\textit{LQCoh}^{fbc}(\mathcal{O}_\mathcal{X})$.
\item Given $\mathcal{F}$ of finite presentation and $\mathcal{G}$ in
$\textit{LQCoh}^{fbc}(\mathcal{O}_\mathcal{X})$ then
$\SheafHom_{\mathcal{O}_\mathcal{X}}(\mathcal{F}, \mathcal{G})$
is in $\textit{LQCoh}^{fbc}(\mathcal{O}_\mathcal{X})$.
\end{enumerate}
\end{proposition}

\begin{proof}
Part (1) is
Sheaves on Stacks, Lemma
\ref{stacks-sheaves-lemma-lqc-flat-base-change-fppf-sheaf}.

\medskip\noindent
Part (2) for the embedding
$\textit{LQCoh}^{fbc}(\mathcal{O}_\mathcal{X}) \subset
\textit{Mod}(\mathcal{X}_\etale, \mathcal{O}_\mathcal{X})$
we have seen in the proof of
Lemma \ref{lemma-flat-comparison}.
Let us prove (2) for the embedding
$\textit{LQCoh}^{fbc}(\mathcal{O}_\mathcal{X}) \subset
\textit{Mod}(\mathcal{O}_\mathcal{X})$.
Let $\varphi : \mathcal{F} \to \mathcal{G}$ be a morphism between
objects of $\textit{LQCoh}^{fbc}(\mathcal{O}_\mathcal{X})$. Since
$\Ker(\varphi)$
is the same whether computed in the \'etale or the fppf
topology, we see that $\Ker(\varphi)$ is in
$\textit{LQCoh}^{fbc}(\mathcal{O}_\mathcal{X})$
by the \'etale case. On the other hand,
the cokernel computed in the fppf topology is the fppf sheafification
of the cokernel computed in the \'etale topology. However, this
\'etale cokernel is in $\textit{LQCoh}^{fbc}(\mathcal{O}_\mathcal{X})$
hence an fppf sheaf
by (1) and we see that the cokernel is in
$\textit{LQCoh}^{fbc}(\mathcal{O}_\mathcal{X})$.
Finally, suppose that
$$
0 \to \mathcal{F}_1 \to \mathcal{F}_2 \to \mathcal{F}_3 \to 0
$$
is an exact sequence in $\textit{Mod}(\mathcal{O}_\mathcal{X})$
(i.e., using the fppf topology) with $\mathcal{F}_1$, $\mathcal{F}_2$
in $\textit{LQCoh}^{fbc}(\mathcal{O}_\mathcal{X})$. In order to show that
$\mathcal{F}_2$ is an object of $\textit{LQCoh}^{fbc}(\mathcal{O}_\mathcal{X})$
it suffices to show that
the sequence is also exact in the \'etale topology. To do this it
suffices to show that any element of $H^1_{fppf}(x, \mathcal{F}_1)$
becomes zero on the members of an \'etale covering of $x$ (for any
object $x$ of $\mathcal{X}$). This is true because
$H^1_{fppf}(x, \mathcal{F}_1) = H^1_\etale(x, \mathcal{F}_1)$ by
Sheaves on Stacks, Lemma \ref{stacks-sheaves-lemma-compare-fppf-etale}
and because of locality of cohomology, see
Cohomology on Sites, Lemma
\ref{sites-cohomology-lemma-kill-cohomology-class-on-covering}.
This proves (2).

\medskip\noindent
Part (3) follows from
Lemma \ref{lemma-check-flat-comparison-on-etale-covering}
and
Sheaves on Stacks, Lemma \ref{stacks-sheaves-lemma-pullback-lqc}.

\medskip\noindent
Part (4)(b) for $R^if_*\mathcal{F}$ computed in the \'etale cohomology
follows from Lemma \ref{lemma-flat-comparison}.
Whereupon part (4)(a) follows from
Sheaves on Stacks, Lemma \ref{stacks-sheaves-lemma-compare-fppf-etale}
combined with (1) above.

\medskip\noindent
Part (5) for the \'etale topology follows from
Sheaves on Stacks, Lemma \ref{stacks-sheaves-lemma-lqc-colimits} and
Lemma \ref{lemma-check-flat-comparison-on-etale-covering}.
The fppf version then follows as the colimit in the \'etale
topology is already an fppf sheaf by part (1).

\medskip\noindent
Parts (6) and (7) follow from the corresponding parts of
Lemma \ref{lemma-check-flat-comparison-on-etale-covering} and
Sheaves on Stacks, Lemma \ref{stacks-sheaves-lemma-lqc-colimits}.
\end{proof}

\end{document}
